% Generated draft: requires mathematical and source review.
\documentclass[11pt]{article}
\usepackage[margin=25mm]{geometry}
\usepackage{fontspec}
\setmainfont{Latin Modern Roman}
\usepackage{amsmath,amssymb,mathtools,unicode-math}
\setmathfont{Latin Modern Math}
\usepackage{enumitem,needspace,xurl,hyperref}
\setlength{\parindent}{0pt}
\setlength{\parskip}{5pt}
\setlength{\emergencystretch}{3em}
\newcommand{\N}{\mathbb{N}}
\newcommand{\Z}{\mathbb{Z}}
\newcommand{\Q}{\mathbb{Q}}
\newcommand{\R}{\mathbb{R}}
\newcommand{\C}{\mathbb{C}}
\newcommand{\F}{\mathbb{F}}
\newcommand{\A}{\mathbb{A}}
\newcommand{\PP}{\mathbb{P}}
\newcommand{\E}{\mathbb{E}}
\newcommand{\eps}{\varepsilon}
\newcommand{\dd}{\,\mathrm{d}}
\begin{document}
\section*{UnsolvedMath ID: 20001152}
\textbf{Source problem code:} catalogue entry
\subsection{Definitions and mathematical statement}
Work with finite simple undirected graphs. A hole is an induced cycle of length at least four; a graph is even-hole-free if it has no induced cycle of even length at least four. For \(v\in V(G)\), let \(N_G(v)\) be its open neighborhood. A vertex is bisimplicial if \(N_G(v)=A\mathbin{\dot\cup}B\), where \(A,B\) each induce a clique; empty parts are allowed. Equivalently, its neighborhood is covered by at most two cliques.

Conjecture 4 of Hayward and Reed states that every nonempty even-hole-free graph has a bisimplicial vertex. The nonempty qualification is necessary for an existential vertex assertion. The supplied research summary reports that Chudnovsky and Seymour proved this statement in 2023, and gives the stronger rooted theorem: if \(K\) is a clique with \(|K|\le2\) and
\[
N_G[K]=K\cup\bigcup_{k\in K}N_G(k)\ne V(G),
\]
then some bisimplicial vertex belongs to \(V(G)\setminus N_G[K]\). Thus ``nondominating'' refers to failure of the closed neighborhood to cover the graph. The convention \(N_G[\varnothing]=\varnothing\) permits the empty clique.

The aid also records the following elimination consequences. Given such a graph \(G\) and any clique \(K\) of size at most two, put \(C=N_G[K]\), \(M=V(G)\setminus C\), and \(t=|M|\). There is an ordering \(v_1,\ldots,v_t\) of \(M\) such that, in
\[
G_i=G-\{v_1,\ldots,v_{i-1}\},
\]
every \(v_i\) is bisimplicial. All vertices of the original core \(C\) remain present throughout this prefix. The statement is vacuous when \(M\) is empty. It gives the exact identity
\[
|E(G)|-|E(G[C])|=\sum_{i=1}^t d_{G_i}(v_i).
\]
Appending a bisimplicial elimination ordering of the induced graph \(G[C]\) gives an ordering of all vertices. Orienting each edge from its earlier endpoint to its later endpoint produces an acyclic orientation in which every out-neighborhood is covered by two cliques. The protected-prefix ordering and edge identity are supplied corollaries, distinct from the original existential conjecture.

The registry records solved status and classification solved in literature; attempt 1 is labeled full solution, proved, with low novelty confidence, and the problem was labeled solved at that run. These status labels do not amount to an independent review of the proof or its corollaries. The accompanying research aid is machine generated, not peer reviewed, and requires independent verification before citation.
\subsection{Short English statement}
Every nonempty graph without an induced even cycle of length at least four is asserted to have a vertex whose neighbors split into at most two cliques. The supplied literature report says this conjecture was proved and gives a stronger version that allows vertices outside a small clique's closed neighborhood to be eliminated first.
\subsection{Sources}
\begin{itemize}
\item Local UnsolvedMath registry, permanent ID \(20001152\), source code \texttt{catalogue entry}; canonical location \texttt{aim-combinatorics-notes.json}, \texttt{notes[276]}.
\item American Institute of Mathematics, \emph{The Perfect Graph Conjecture}, Conjecture 4 (Hayward and Reed): \url{https://aimath.org/WWN/perfectgraph/perfectgraph.pdf}. The original literature field is empty.
\item Chudnovsky and Seymour, \emph{Journal of Combinatorial Theory, Series B} \textbf{161} (2023), 331--381, as cited in the supplied research summary for the bisimplicial-vertex and rooted theorems. An article title and theorem number are not supplied.
\item Supplied research attempt 1 and registry update of August 14, 2026; machine-generated aid, not peer reviewed.
\end{itemize}
\end{document}
